\chapter{Fichas de las capas}
\label{cap:anexol}

\section*{L.1. Alcance}

Este anexo resume la finalidad, entradas, salidas, evaluación y limitaciones de las
tres capas. La separación de responsabilidades es esencial: Capa 1 estructura el texto,
Capa 2 recomienda la prioridad y Capa 3 explica esa recomendación. La persona operadora
mantiene la decisión final.

Los detalles anteriormente distribuidos entre los anexos técnicos de recuperación
normativa y configuración del LLM se consolidan aquí para ofrecer una visión académica
única y evitar inventarios de rutas, clases o ficheros.

\section*{L.2. Capa 1: extracción NLP simbólica}

\textbf{Finalidad.} Transformar el título y la descripción del incidente en variables
predecisionales estructuradas. El procedimiento combina detección semántica mediante
reglas y diccionarios operativos, tratamiento de negaciones, normalización de valores y
estimación de confianza. La salida conserva la señal activada y su procedencia para
facilitar la auditoría.

\textbf{Evaluación.} Sobre cinco escenarios curados obtiene precisión micro
(0{,}733333), sensibilidad micro (1{,}000000), F1 micro (0{,}846154) y latencia
media (4{,}58844) ms. Es una comprobación funcional reducida, no una validación sobre
llamadas reales ni sobre un conjunto de referencia etiquetado por el 112.

\textbf{Limitación.} La versión evaluada no contiene un clasificador aprendido
congelado ni se le atribuyen entrenamientos o métricas procedentes de la Capa 2. La
posible incorporación de modelos aprendidos queda como trabajo futuro y requeriría una
evaluación específica.

\section*{L.3. Capa 2: priorización interpretable}

\subsection*{L.3.1. Datos y etiqueta objetivo}

La Capa 2 se entrena y evalúa con 9.380 registros limpios del Registro 112 Castilla y
León. La prioridad P1--P4 no existe en la fuente: es una etiqueta académica construida
mediante supervisión débil. El acuerdo principal alcanza un alfa de Krippendorff de
\(\alpha = 0{,}899902\). Las particiones utilizadas exclusivamente para Capa 2 son 6.512
registros de entrenamiento, 1.415 de validación y 1.453 de prueba; la prueba temporal de
2022 contiene 735 registros.

Las variables de entrada se limitan a señales disponibles antes de la primera decisión:
fallecimiento, herido grave, atrapamiento, intoxicación, llamadas múltiples, incendio,
accidente de tráfico, rescate, meteorología o inundación, riesgo vital textual y
categoría preliminar normalizada. Se excluyen recursos movilizados, pacientes atendidos,
cierre del incidente y cualquier información posterior, con el fin de evitar fugas de
información.

\subsection*{L.3.2. Alternativas y selección}

Se comparan una línea base de 19 reglas expertas, RuleFit-lite y una implementación
externa de RuleFit mediante \texttt{imodels} \cite{Singh2021imodels}. Esta última se
utiliza con finalidad diagnóstica. La selección se realiza con la partición de validación
y la estimación final se obtiene una sola vez sobre la partición de prueba.

\begin{table}[htbp]
\centering
\caption{Resultados de los modelos de Capa 2 en la partición de prueba}
\label{tab:model-card-capa2-comparativa}
\small
\begin{tabular}{|l|r|r|r|r|}
\hline
\textbf{Modelo} & \textbf{Exactitud} & \textbf{Macro-F1} & \textbf{Sensibilidad P1} & \textbf{Reglas} \\ \hline
Línea base heurística & 0,694425 & 0,464284 & 0,998473 & 19 \\ \hline
RuleFit-lite & 0,903648 & 0,905604 & 0,909924 & 30 \\ \hline
RuleFit externo & 0,768385 & 0,568239 & 0,902778 & 30 \\ \hline
\end{tabular}
\end{table}

RuleFit-lite se selecciona porque presenta la mejor macro-F1, mantiene una sensibilidad
P1 superior al umbral interno de (0{,}85), conserva reglas exportables y ofrece menor
coste local que la configuración externa ensayada. En la prueba temporal alcanza
exactitud (0{,}884354), macro-F1 (0{,}877594) y sensibilidad P1 (0{,}928767).

Estas métricas corresponden al estimador RuleFit-lite \textbf{sin Puerta de Seguridad}.
La salvaguarda es una función determinista de posprocesamiento aplicada durante la
inferencia integrada; se comprueba mediante escenarios funcionales y no se mezcla con
la comparación estadística de modelos.

\subsection*{L.3.3. Riesgos y trazabilidad}

La matriz de confusión identifica 59 falsos negativos P1: 34 casos clasificados como
P2, 18 como P3 y 7 como P4. Los siete errores P1--P4 constituyen el riesgo principal y
forman parte de la muestra de 119 casos preparada para revisión. Los campos reservados
a la persona revisora permanecen vacíos, por lo que esta muestra no equivale a una
validación humana completada.

\section*{L.4. Capa 3: explicación local y recuperación normativa}

\textbf{Finalidad.} Capa 3 recibe la prioridad, las probabilidades, las reglas activadas
y el contexto explicativo generado por Capa 2. Produce una explicación estructurada y
evidencias normativas, pero no recalcula ni modifica la prioridad.

\textbf{Recuperación.} El sistema utiliza un índice semántico local sobre normativa
estatal y de Castilla y León. Los documentos se segmentan y representan con el modelo
multilingüe \texttt{paraphrase-multilingual-MiniLM-L12-v2}. La cobertura incluye la
legislación de protección civil y los planes territoriales y especiales disponibles. No
se dispone del texto técnico completo de MPCyL, por lo que la recuperación en escenarios
químicos debe interpretarse con cautela y completarse en una versión futura con la
fuente oficial íntegra.

\textbf{Generación y salvaguardas.} El LLM se ejecuta localmente con temperatura
(0{,}0). Las instrucciones exigen mantener la prioridad de Capa 2, utilizar solo las
reglas y evidencias recibidas, declarar incertidumbre y evitar hechos o citas no
recuperados. Las herramientas de consulta normativa y detalle de reglas se exponen de
forma controlada mediante Model Context Protocol. Si el modelo no responde o incumple
el contrato, se genera una explicación alternativa determinista que conserva la salida
de Capa 2.

\textbf{Evaluación.} La prueba funcional con LLM local resuelve correctamente cuatro
escenarios P1--P4 y observa una latencia media de (4{,}898) s en el entorno
experimental. La evaluación automatizada de fidelidad sobre los 119 casos preparados
alcanza una tasa de cumplimiento de (1{,}000000), fidelidad media (0{,}957367) y
mínima (0{,}933333). Se trata de un juez determinista sobre explicaciones controladas,
no de una evaluación humana ni de un LLM-as-Judge independiente.

\section*{L.5. Entrada de voz experimental}

La interfaz incorpora reconocimiento automático del habla mediante
\texttt{faster-whisper}. Su salida propone texto y campos normalizados que deben
revisarse antes de ejecutar la inferencia. Esta función se ha utilizado como demostración
con audio de prueba; no se reportan tasas de error de palabra ni validación con llamadas
reales, ruido de sala o diversidad de locutores.

\section*{L.6. Limitaciones comunes}

\begin{itemize}
    \item La referencia P1--P4 es académica y débil, no una prioridad oficial del 112.
    \item Las pruebas integradas muestran compatibilidad técnica en escenarios
    controlados, no utilidad operativa ni cumplimiento de un acuerdo de nivel de servicio.
    \item La Puerta de Seguridad es una salvaguarda conservadora no certificada y debe
    evaluarse separadamente del modelo estadístico.
    \item La muestra de 119 casos está preparada para revisión, pero no contiene
    decisiones humanas completadas.
    \item Los resultados no acreditan conformidad regulatoria ni autorizan el uso del
    prototipo en un servicio de emergencias real.
\end{itemize}
